\section{The Appointed Texts}
\label{sec:appointed-texts}

The Latin below is that of the 1962 typical edition, \work{Missale Romanum}
(Typis Polyglottis Vaticanis), printed pp.~421--422, nos.~1759--1768, read on
the page images of the CMAA facsimile. The Missal prints the three orations,
the Epistle and the Gospel a second time in the season after the Epiphany
(pp.~48--49, nos.~337--345), with the same words. The Latin here follows the
resumed printing: the three orations close with the cue \latin{Per Dóminum.},
and the lessons keep that printing's pointing. Every word of all ten texts
stands in the Ratisbon printing of the \work{Missale Romanum} (Pustet, 1862):
the orations and lessons under the Sixth Sunday after the Epiphany,
pp.~44--45, and the chants under the Twenty-third Sunday after Pentecost,
pp.~357--358, where a rubric directs that the Sundays resumed after Pentecost
take those chants. The two printings differ in spelling (\latin{exequámur},
\latin{expectáre}, \latin{charitátis}, \latin{cœlis}), in pointing and in the
form of the citations, and in no word that is said.

The Lasance \work{New Roman Missal}, in the Benziger revision of 1945, prints
this whole Mass in Latin and English as its \work{Supplementary Mass No.~1},
``Mass of the Sixth Sunday after Epiphany,'' pp.~740--743, under the rubric
that it is said on the twenty-fifth Sunday after Pentecost when there are
twenty-six. Its English of the five chants renders the Missal's own forms of
the verses where they differ from the Clementine Vulgate, and that English is
given for each chant; its English of the three orations is given for them. The
Douay--Rheims Bible in Bishop Challoner's revision, which translates the
Clementine Vulgate, is given for the two lessons and beside each chant at the
canonical verses (for the Offertory, which repeats the Alleluia's verses,
under the Alleluia), and each place where the Missal's form departs from the
Vulgate is named beneath the text. None of this English is an approved
liturgical translation of the 1962 Missal. Psalms are cited in the Missal's
numbering throughout; the Hebrew numbering is added in parentheses at each
appointed psalm text below and in the map of the texts.

\subsection{Introit}
\label{proper-introit}

\properlocus{\latin{Antiphona ad Introitum}, \latin{Ierem.\ 29, 11, 12 et 14};
verse \latin{Ps.\ 84, 2} (Heb.\ 85). Missal p.~421, no.~1759.}

\begin{latinproper}
Dicit Dóminus: Ego cógito cogitatiónes pacis, et non afflictiónis:
invocábitis me, et ego exáudiam vos: et redúcam captivitátem vestram de
cunctis locis. \textnormal{Ps.\ 84, 2} Benedixísti, Dómine, terram tuam:
avertísti captivitátem Iacob. ℣.~Glória Patri.
\end{latinproper}

\begin{englishwitness}{Lasance, The New Roman Missal (1945), Supplementary Mass No.~1, Introit}
The Lord saith, I think thoughts of peace, and not of affliction; you shall
call upon Me and I will hear you; and I will bring back your captivity from
all places. Ps.~84, 2. Lord, Thou hast blessed Thy land: Thou hast turned away
the captivity of Jacob. ℣.~Glory.
\end{englishwitness}

\begin{englishwitness}{Douay--Rheims (Challoner), Jer 29:11--12, 14, the antiphon's verses}
For I know the thoughts that I think towards you, saith the Lord, thoughts of
peace, and not of affliction, to give you an end and patience. And you shall
call upon me, and you shall go: and you shall pray to me, and I will hear you.
\dots\ And I will be found by you, saith the Lord: and I will bring back your
captivity, and I will gather you out of all nations, and from all the places
to which I have driven you out, saith the Lord \dots
\end{englishwitness}

\noindent The antiphon is not a continuous verse of the Vulgate. It opens
\latin{Dicit Dóminus}, where the Vulgate's verse~11 has God say \latin{Ego
enim scio cogitationes quas ego cogito super vos, ait Dominus}. It has nothing
of \latin{ut dem vobis finem et patientiam} (``to give you an end and
patience''), of \latin{et ibitis: et orabitis me} (``and you shall go: and you
shall pray to me''), of verse~13, or of \latin{et congregabo vos de universis
gentibus} (``and I will gather you out of all nations''), and it joins
\latin{redúcam captivitátem vestram} directly to \latin{de cunctis locis}. The
psalm verse agrees with the Vulgate. After the \latin{Gloria Patri} the
antiphon is repeated.

\subsection{Collect}
\label{proper-collect}

\properlocus{\latin{Oratio}. Missal p.~421, no.~1760.}

\begin{latinproper}
Præsta, quǽsumus, omnípotens Deus: ut, semper rationabília meditántes, quæ
tibi sunt plácita, et dictis exsequámur et factis. Per Dóminum.
\end{latinproper}

\begin{englishwitness}{Lasance, The New Roman Missal (1945), Supplementary Mass No.~1, Prayer}
Grant, we beseech Thee, almighty God, that, ever fixing our thoughts on
reasonable things, we may both in word and in deed do what is pleasing to
Thee. Through our Lord.
\end{englishwitness}

\noindent The Epiphany-season printing of the same prayer (no.~337) writes the
conclusion out and sets a comma before \latin{et factis}. Whichever cue a
printing gives, the prayer is said with the Missal's one fixed conclusion.

\subsection{Epistle}
\label{proper-epistle}

\properlocus{\latin{Léctio Epístolæ beáti Pauli Apóstoli ad Thessalonicénses},
1~Thess 1:2--10. Missal pp.~421--422, no.~1761.}

\begin{latinproper}
Fratres: Grátias ágimus Deo semper pro ómnibus vobis, memóriam vestri
faciéntes in oratiónibus nostris sine intermissióne, mémores óperis fidei
vestræ, et labóris, et caritátis, et sustinéntiæ spei Dómini nostri Iesu
Christi, ante Deum et Patrem nostrum: sciéntes, fratres, dilécti a Deo,
electiónem vestram: quia Evangélium nostrum non fuit ad vos in sermóne
tantum, sed et in virtúte, et in Spíritu Sancto, et in plenitúdine multa,
sicut scitis quales fuérimus in vobis propter vos. Et vos imitatóres nostri
facti estis, et Dómini, excipiéntes verbum in tribulatióne multa, cum gáudio
Spíritus Sancti: ita ut facti sitis forma ómnibus credéntibus in Macedónia,
et in Acháia. A vobis enim diffamátus est sermo Dómini, non solum in
Macedónia, et in Acháia, sed et in omni loco fides vestra, quæ est ad Deum,
profécta est, ita ut non sit nobis necésse quidquam loqui. Ipsi enim de nobis
annúntiant qualem intróitum habuérimus ad vos: et quómodo convérsi estis ad
Deum a simulácris, servíre Deo vivo et vero, et exspectáre Fílium eius de
cælis (quem suscitávit ex mórtuis) Iesum, qui erípuit nos ab ira ventúra.
\end{latinproper}

\begin{englishwitness}{Douay--Rheims (Challoner), 1 Thess 1:2--10}
Grace be to you and peace. We give thanks to God always for you all: making a
remembrance of you in our prayers without ceasing, Being mindful of the work
of your faith and labour and charity: and of the enduring of the hope of our
Lord Jesus Christ before God and our Father. Knowing, brethren, beloved of
God, your election: For our gospel hath not been unto you in word only, but in
power also: and in the Holy Ghost and in much fulness, as you know what manner
of men we have been among you for your sakes. And you became followers of us
and of the Lord: receiving the word in much tribulation, with joy of the Holy
Ghost: So that you were made a pattern to all that believe in Macedonia and in
Achaia. For from you was spread abroad the word of the Lord not only in
Macedonia and in Achaia but also in every place: your faith which is towards
God, is gone forth, so that we need not to speak any thing. For they
themselves relate of us, what manner of entering in we had unto you: and how
you turned to God from idols to serve the living and true God. And to wait for
his Son from heaven (whom he raised up from the dead), Jesus, who hath
delivered us from the wrath to come.
\end{englishwitness}

\noindent The lesson begins within the Vulgate's verse~2: the Missal sets
\latin{Fratres} where the Vulgate opens the verse with \latin{Gratia vobis, et
pax} (``Grace be to you and peace''), and in verse~10 it reads \latin{ex
mórtuis} where the Vulgate reads \latin{a mortuis}, a difference the English
does not show. The resumed printing heads the lesson \latin{1 Thess.\ 1,
1-20}; the chapter has ten verses, and the same lesson in the Epiphany season
is headed \latin{1 Thess.\ 1, 2-10}, the extent of the text printed.

\subsection{Gradual}
\label{proper-gradual}

\properlocus{\latin{Graduale}, \latin{Ps.\ 43, 8-9} (Heb.\ 44). Missal p.~422,
no.~1762.}

\begin{latinproper}
Liberásti nos, Dómine, ex affligéntibus nos: et eos, qui nos odérunt,
confudísti. ℣.~In Deo laudábimur tota die, et in nómine tuo confitébimur in
sǽcula.
\end{latinproper}

\begin{englishwitness}{Lasance, The New Roman Missal (1945), Supplementary Mass No.~1, Gradual}
Thou hast delivered us, O Lord, from them that afflict us: and hast put them
to shame that hate us. ℣.~In God we will glory all the day: and in Thy name we
will give praise for ever.
\end{englishwitness}

\begin{englishwitness}{Douay--Rheims (Challoner), Ps 43:8--9}
But thou hast saved us from them that afflict us: and hast put them to shame
that hate us. In God shall we glory all the day long: and in thy name we will
give praise for ever.
\end{englishwitness}

\noindent The Vulgate's verse~8 reads \latin{salvasti enim nos de affligentibus
nos, et odientes nos confudisti}; the Gradual has \latin{Liberásti nos, Dómine,
ex affligéntibus nos: et eos, qui nos odérunt, confudísti}. Verse~9 ends
\latin{in saeculum} in the Vulgate and \latin{in sǽcula} in the Gradual.

\subsection{Alleluia}
\label{proper-alleluia}

\properlocus{\latin{Ps.\ 129, 1-2} (Heb.\ 130). Missal p.~422, no.~1763.}

\begin{latinproper}
Allelúia, allelúia. ℣.~De profúndis clamávi ad te, Dómine: Dómine, exáudi
oratiónem meam. Allelúia.
\end{latinproper}

\begin{englishwitness}{Lasance, The New Roman Missal (1945), Supplementary Mass No.~1, Alleluia}
℣.~From the depths I have cried to Thee, O Lord: Lord, hear my prayer.
\end{englishwitness}

\begin{englishwitness}{Douay--Rheims (Challoner), Ps 129:1--2}
Out of the depths I have cried to thee, O Lord: Lord, hear my voice. Let thy
ears be attentive to the voice of my supplication.
\end{englishwitness}

\noindent The Vulgate reads \latin{Domine, exaudi vocem meam}, ``Lord, hear my
voice''; the chant reads \latin{oratiónem meam}, ``my prayer,'' and stops
before the second half of verse~2. The Mass has no Tract and no Sequence.

\subsection{Gospel}
\label{proper-gospel}

\properlocus{\latin{Sequéntia sancti Evangélii secúndum Matthǽum}, Mt
13:31--35. Missal p.~422, no.~1764; the printed cue \latin{Credo} follows.}

\begin{latinproper}
In illo témpore: Dixit Iesus turbis parábolam hanc: Símile est regnum cælórum
grano sinápis, quod accípiens homo seminávit in agro suo: quod mínimum quidem
est ómnibus semínibus: cum autem créverit, maius est ómnibus oléribus, et fit
arbor, ita ut vólucres cæli véniant et hábitent in ramis eius. Aliam parábolam
locútus est eis: Símile est regnum cælórum ferménto, quod accéptum múlier
abscóndit in farínæ satis tribus, donec fermentátum est totum. Hæc ómnia
locútus est Iesus in parábolis ad turbas: et sine parábolis non loquebátur
eis: ut implerétur quod dictum erat per Prophétam dicéntem: Apériam in
parábolis os meum, eructábo abscóndita a constitutióne mundi.
\end{latinproper}

\begin{englishwitness}{Douay--Rheims (Challoner), Mt 13:31--35}
Another parable he proposed unto them, saying: The kingdom of heaven is like
to a grain of mustard seed, which a man took and sowed in his field. Which is
the least indeed of all seeds; but when it is grown up, it is greater than all
herbs, and becometh a tree, so that the birds of the air come, and dwell in the
branches thereof. Another parable he spoke to them: The kingdom of heaven is
like to leaven, which a woman took and hid in three measures of meal, until
the whole was leavened. All these things Jesus spoke in parables to the
multitudes: and without parables he did not speak to them. That it might be
fulfilled which was spoken by the prophet, saying: I will open my mouth in
parables, I will utter things hidden from the foundation of the world.
\end{englishwitness}

\noindent The Missal's opening, \latin{In illo témpore: Dixit Iesus turbis
parábolam hanc}, which the Lasance renders ``At that time, Jesus spoke this
parable to the multitudes,'' stands for the Vulgate's \latin{Aliam parabolam
proposuit eis dicens}. Every other word of verses 31--35 agrees with the
Vulgate.

\subsection{Offertory}
\label{proper-offertory}

\properlocus{\latin{Antiphona ad Offertorium}, \latin{Ps.\ 129, 1-2} (Heb.\ 130).
Missal p.~422, no.~1765.}

\begin{latinproper}
De profúndis clamávi ad te, Dómine: Dómine, exáudi oratiónem meam: de
profúndis clamávi ad te, Dómine.
\end{latinproper}

\begin{englishwitness}{Lasance, The New Roman Missal (1945), Supplementary Mass No.~1, Offertory}
From the depths I have cried out to Thee, O Lord; Lord, hear my prayer: from
the depths I have cried out to Thee, O Lord.
\end{englishwitness}

\noindent The Offertory sings the Alleluia's verses again, with the same
\latin{oratiónem meam}, and repeats its first line at the end.

\subsection{Secret}
\label{proper-secret}

\properlocus{\latin{Secreta}. Missal p.~422, no.~1766; the printed direction
\latin{Præfatio de Ssma Trinitate} follows.}

\begin{latinproper}
Hæc nos oblátio, Deus, mundet, quǽsumus, et rénovet, gubérnet et prótegat.
Per Dóminum.
\end{latinproper}

\begin{englishwitness}{Lasance, The New Roman Missal (1945), Supplementary Mass No.~1, Secret}
May this oblation, O God, cleanse, renew, govern, and protect us, we beseech
Thee. Through our Lord.
\end{englishwitness}

\subsection{Communion}
\label{proper-communion}

\properlocus{\latin{Antiphona ad Communionem}, \latin{Marc.\ 11, 24}. Missal
p.~422, no.~1767.}

\begin{latinproper}
Amen dico vobis, quidquid orántes pétitis, crédite quia accipiétis, et fiet
vobis.
\end{latinproper}

\begin{englishwitness}{Lasance, The New Roman Missal (1945), Supplementary Mass No.~1, Communion}
Amen I say to you, whatsoever you ask when you pray, believe that you shall
receive, and it shall be done to you.
\end{englishwitness}

\begin{englishwitness}{Douay--Rheims (Challoner), Mk 11:23--24}
Amen I say to you that whosoever shall say to this mountain, Be thou removed
and be cast into the sea, and shall not stagger in his heart, but believe that
whatsoever he saith shall be done; it shall be done unto him. Therefore I say
unto you, all things, whatsoever you ask when ye pray, believe that you shall
receive: and they shall come unto you.
\end{englishwitness}

\noindent The Vulgate's verse~24 reads \latin{Propterea dico vobis, omnia
quaecumque orantes petitis, credite quia accipietis, et evenient vobis}. The
antiphon differs at both ends: it opens \latin{Amen dico vobis} and closes
\latin{et fiet vobis}. In Mark the verse before it, the saying about the
mountain, also opens \latin{Amen dico vobis} and closes \latin{fiet ei}. The
Lasance heads the antiphon ``Mark 11, 34,'' a misprint for verse~24.

\subsection{Postcommunion}
\label{proper-postcommunion}

\properlocus{\latin{Postcommunio}. Missal p.~422, no.~1768.}

\begin{latinproper}
Cæléstibus, Dómine, pasti delíciis: quǽsumus; ut semper éadem, per quæ
veráciter vívimus, appetámus. Per Dóminum.
\end{latinproper}

\begin{englishwitness}{Lasance, The New Roman Missal (1945), Supplementary Mass No.~1, Postcommunion}
Being fed with celestial delights, we beseech Thee, O Lord, that we may ever
hunger after those things by which we truly live. Through our Lord.
\end{englishwitness}
